\documentclass[12pt]{article}
\usepackage[utf8]{inputenc}
\usepackage[T1]{fontenc}
\usepackage[margin=1in]{geometry}
\usepackage{setspace}
\usepackage{xcolor}
\usepackage[colorlinks=true, urlcolor=blue!60!black, linkcolor=blue!60!black]{hyperref}
\setlength{\parindent}{2em}
\setlength{\parskip}{1.25em}
\renewcommand{\baselinestretch}{1.0}
\begin{document}
\begin{center}
{\LARGE \textit{San Antonio Independent School District v. Rodriguez}, 411 U.S. 1 (1973)\par}
\vspace{0.8em}
{\large Case Brief\par}
\end{center}
\vspace{1em}

\section{Memory Jogger}

In a 5--4 decision marking the end of the Warren Court's expansion of Equal Protection doctrine, the Supreme Court held that Texas's property-tax-based school funding system did not violate the Fourteenth Amendment because neither poverty nor district wealth is a suspect classification and education is not a fundamental federal constitutional right.

\section{Detailed Case Facts}

Texas financed its public elementary and secondary schools through a dual state-and-local system. In 1970--1971, roughly 48\% of total public school funding came from the state's Minimum Foundation School Program (created by the 1949 Gilmer-Aikin bills), 41.1\% came from local ad valorem property taxes, and 10.9\% came from federal funds (\href{https://supreme.justia.com/cases/federal/us/411/1/}{Justia}). Under the Minimum Foundation Program---designed to guarantee a basic minimum education in every district---the State funded approximately 80\% of the program from general revenues, while Texas's 1,161 local school districts collectively funded the remaining 20\% through a ``Local Fund Assignment'' apportioned by a county economic index (\href{https://supreme.justia.com/cases/federal/us/411/1/}{Justia}). Every school district then supplemented its state allotment by levying local property taxes on assessable property within its borders, subject to a statutory maintenance tax ceiling of \$1.50 per \$100 of valuation (Tex. Educ. Code Ann. \S~20.04(d)) (\href{https://supreme.justia.com/cases/federal/us/411/1/}{Justia}).

Because commercial, industrial, and residential property values varied dramatically across districts, this reliance on local property taxes produced stark inter-district spending disparities. Throughout the litigation, the parties compared two San Antonio metropolitan districts for the 1967--1968 school year (\href{https://supreme.justia.com/cases/federal/us/411/1/}{Justia}):
\begin{itemize}
\item \textbf{Edgewood Independent School District:} Enrolled approximately 22,000 students in 25 schools in a core-city residential sector with little commercial property; roughly 90\% of students were Mexican-American and over 6\% were Black (\href{https://supreme.justia.com/cases/federal/us/411/1/}{Justia}). Edgewood had the lowest average assessed property value per pupil in the metropolitan area (\$5,960) and the lowest median family income (\$4,686 under 1960 census data) (\href{https://supreme.justia.com/cases/federal/us/411/1/}{Justia}). Even at the area's highest equalized tax rate (\$1.05 per \$100 of assessed property), Edgewood raised only \$26 per pupil above its Local Fund Assignment; combined with \$222 from the Foundation Program and \$108 in federal aid, Edgewood spent \textbf{\$356 per pupil} (\href{https://supreme.justia.com/cases/federal/us/411/1/}{Justia}).
\item \textbf{Alamo Heights Independent School District:} Enrolled approximately 5,000 students in 6 schools in an affluent, predominantly ``Anglo'' residential community (18\% Mexican-American and under 1\% Black) (\href{https://supreme.justia.com/cases/federal/us/411/1/}{Justia}). Its assessed property value exceeded \$49,000 per pupil (\$49,078 in Justice White's dissent), and its median family income was \$8,001 (\href{https://supreme.justia.com/cases/federal/us/411/1/}{Justia}). At a lower equalized tax rate (\$0.85 per \$100), Alamo Heights raised \$333 per pupil locally above its Foundation contribution; combined with \$225 from the Foundation Program and \$36 in federal aid, Alamo Heights spent \textbf{\$594 per pupil} (\href{https://supreme.justia.com/cases/federal/us/411/1/}{Justia}).
\end{itemize}

Although state Foundation aid increased substantially by 1970--1971 (to \$356 per pupil in Edgewood and \$491 in Alamo Heights), Alamo Heights received more state aid per pupil because the Foundation formula paid higher minimum salaries to teachers with more experience and advanced degrees (\href{https://supreme.justia.com/cases/federal/us/411/1/}{Justia}). Demetrio P. Rodriguez and other Mexican-American parents of Edgewood schoolchildren filed a class action on behalf of schoolchildren throughout Texas who are members of minority groups or poor and reside in school districts with low property tax bases (\href{https://supreme.justia.com/cases/federal/us/411/1/}{Justia}).

\section{Procedural History}

\begin{itemize}
\item \textbf{Filing and Pretrial Dismissal of School Districts (1968--1969):} Plaintiffs filed suit in the summer of 1968 (June 30, 1968) in the U.S. District Court for the Western District of Texas against the State Board of Education, the Commissioner of Education, the State Attorney General, the Bexar County Board of Trustees, and seven San Antonio-area school districts (\href{https://supreme.justia.com/cases/federal/us/411/1/}{Justia}; \href{https://en.wikipedia.org/wiki/San_Antonio_Independent_School_District_v._Rodriguez}{Wikipedia}). After a pretrial conference, the district court dismissed the seven school districts as defendants; San Antonio Independent School District (whose name remained on the caption) later joined plaintiffs' side as amicus curiae in the Supreme Court (\href{https://supreme.justia.com/cases/federal/us/411/1/}{Justia}).
\item \textbf{Three-Judge District Court (1969--1971):} A three-judge federal district court was impaneled in January 1969 under 28 U.S.C. \S\S~1253 and 2281, and trial was delayed two years for discovery and a pending Texas legislative study (\href{https://supreme.justia.com/cases/federal/us/411/1/}{Justia}). In December 1971, the panel unanimously held per curiam that wealth is a suspect classification and education is a fundamental interest, and that the Texas school finance system failed both strict scrutiny and rational basis review (\textit{Rodriguez v. San Antonio Independent School District}, 337 F. Supp. 280 (W.D. Tex. 1971)) (\href{https://supreme.justia.com/cases/federal/us/411/1/}{Justia}). The court stayed its mandate for two years to give Texas time to enact a fiscally neutral plan (\href{https://supreme.justia.com/cases/federal/us/411/1/}{Justia}).
\item \textbf{U.S. Supreme Court (1972--1973):} The State took a direct appeal under 28 U.S.C. \S~1253. The Supreme Court noted probable jurisdiction, 406 U.S. 966 (1972), heard oral argument on October 12, 1972, and reversed on March 21, 1973 (\href{https://www.courtlistener.com/opinion/108751/san-antonio-independent-school-district-v-rodriguez/}{CourtListener}; \href{https://supreme.justia.com/cases/federal/us/411/1/}{Justia}). Rehearing was denied, 411 U.S. 959 (1973).
\end{itemize}

\section{Judicial Votes}

The Supreme Court reversed the three-judge District Court by a \textbf{5--4} vote (\href{https://supreme.justia.com/cases/federal/us/411/1/}{Justia}):
\begin{itemize}
\item \textbf{Majority (5):} Justice Lewis F. Powell, Jr. authored the opinion of the Court, joined by Chief Justice Warren E. Burger and Justices Potter Stewart, Harry A. Blackmun, and William H. Rehnquist.
\item \textbf{Concurrence (1):} Justice Potter Stewart filed a concurring opinion (while also joining the majority opinion in full).
\item \textbf{Dissents (4):}
  \begin{itemize}
  \item Justice William J. Brennan, Jr. filed a solo dissenting opinion (and also joined Justice White's dissent).
  \item Justice Byron R. White filed a dissenting opinion, joined by Justices William O. Douglas and William J. Brennan, Jr.
  \item Justice Thurgood Marshall filed a dissenting opinion, joined by Justice William O. Douglas.
  \item (Justice William O. Douglas did not write separately, joining both Justice White's and Justice Marshall's dissents.)
  \end{itemize}
\end{itemize}

\section{Holding}

\textbf{Judgment:} Reversed. The Texas public school financing system does not violate the Equal Protection Clause of the Fourteenth Amendment (\href{https://supreme.justia.com/cases/federal/us/411/1/}{Justia}).

\textbf{Rule Announced:} Strict scrutiny under the Equal Protection Clause is reserved for state laws that disadvantage a ``suspect class'' (one saddled with disabilities, a history of purposeful unequal treatment, or political powerlessness) or impinge upon a ``fundamental right'' explicitly or implicitly guaranteed by the U.S. Constitution (\href{https://supreme.justia.com/cases/federal/us/411/1/}{Justia}; \href{https://www.law.cornell.edu/constitution-conan/amendment-14/section-1/educational-opportunity}{Constitution Annotated}). Wealth-based disparities across school districts do not define a suspect class where there is no absolute deprivation of education and no definable class of indigent persons, and education is not a fundamental federal constitutional right (\href{https://supreme.justia.com/cases/federal/us/411/1/}{Justia}). Under rational basis review, a state school finance system that guarantees a minimum basic education while relying on local property taxes to promote local participation and control rationally furthers a legitimate state purpose (\href{https://supreme.justia.com/cases/federal/us/411/1/}{Justia}).

\textbf{Holding vs. Dicta:} The \emph{holding} is that relative inter-district disparities in public school funding resulting from local property taxation trigger only rational basis review under the Equal Protection Clause and satisfy that deferential standard (\href{https://supreme.justia.com/cases/federal/us/411/1/}{Justia}). Crucially, the Court's suggestion in footnotes and Part II.B that an \emph{absolute deprivation} of all public education---such as charging tuition that totally excludes indigent children, or failing to provide ``an opportunity to acquire the basic minimal skills necessary'' for speech and voting---might present a ``far more compelling'' constitutional claim is \emph{dicta}, leaving open the door later tested in \href{https://www.law.cornell.edu/supremecourt/text/457/202}{\textit{Plyler v. Doe}}, 457 U.S. 202 (1982) (\href{https://supreme.justia.com/cases/federal/us/411/1/}{Justia}). Likewise, Powell's Part IV ``cautionary postscript'' questioning whether judicially mandated equalization would benefit urban minorities and calling for legislative tax reform is \emph{dicta} (\href{https://supreme.justia.com/cases/federal/us/411/1/}{Justia}).

\section{Analysis of Opinions}

\subsection{Majority Opinion (Justice Powell)}

Justice Powell first rejected the District Court's holding that the Texas system discriminated against a suspect class based on wealth (\href{https://supreme.justia.com/cases/federal/us/411/1/}{Justia}). Examining prior wealth-discrimination cases---such as \textit{Griffin v. Illinois}, 351 U.S. 12 (1956), \textit{Douglas v. California}, 372 U.S. 353 (1963), and \textit{Bullock v. Carter}, 405 U.S. 134 (1972)---Powell distilled two essential features absent here: (1) a definable class of persons completely unable to pay due to indigency, and (2) an \emph{absolute deprivation} of a meaningful opportunity to enjoy the benefit (\href{https://supreme.justia.com/cases/federal/us/411/1/}{Justia}). Here, Texas guaranteed twelve years of free public schooling through its Minimum Foundation Program, and plaintiffs' own statistical evidence (Professor Joel S. Berke's 110-district survey) failed to show that poor families consistently lived in property-poor districts: across the middle 96 districts (nearly 90\% of the sample), the correlation between median family income and educational spending was actually inverted (\href{https://supreme.justia.com/cases/federal/us/411/1/}{Justia}). Nor did residents of property-poor districts constitute a suspect class, as this ``large, diverse, and amorphous class'' lacked the traditional indicia of suspectness (\href{https://supreme.justia.com/cases/federal/us/411/1/}{Justia}).

Second, Powell rejected the claim that education is a fundamental right. While reaffirming \href{https://supreme.justia.com/cases/federal/us/347/483/}{\textit{Brown v. Board of Education}}'s recognition of education's vital societal role, Powell insisted that ``the importance of a service performed by the State does not determine whether it must be regarded as fundamental'' (\href{https://supreme.justia.com/cases/federal/us/411/1/}{Justia}). Citing \textit{Lindsey v. Normet}, 405 U.S. 56 (1972) (housing), and \textit{Dandridge v. Williams}, 397 U.S. 471 (1970) (welfare), he held that the sole test is whether a right is ``explicitly or implicitly guaranteed by the Constitution'' (\href{https://supreme.justia.com/cases/federal/us/411/1/}{Justia}). Education is neither. As for appellees' nexus theory---that education is prerequisite to First Amendment speech and intelligent voting---Powell replied that the Constitution does not guarantee ``the most effective speech or the most informed electoral choice,'' that Texas was not depriving children of basic minimal skills, and that the nexus theory lacked logical limits because food and shelter are equally essential to civic participation (\href{https://supreme.justia.com/cases/federal/us/411/1/}{Justia}). Powell added that judicial restraint was reinforced by the Court's lack of expertise in taxation and educational policy (noting scholarly disagreement over whether spending determines educational quality) and by federalism concerns against invalidating school finance systems in nearly every state (\href{https://supreme.justia.com/cases/federal/us/411/1/}{Justia}).

Third, applying rational basis review, Powell held that Texas's dual system rationally furthered a legitimate state interest: guaranteeing a basic statewide educational floor while preserving local control, fiscal participation, and pluralism in school governance (\href{https://supreme.justia.com/cases/federal/us/411/1/}{Justia}). That some districts enjoyed less fiscal flexibility did not render the system irrational, because rational basis review tolerates ``some inequality'' and does not require the least restrictive alternative (\href{https://supreme.justia.com/cases/federal/us/411/1/}{Justia}).

\subsection{Concurring Opinion (Justice Stewart)}

Justice Stewart joined Powell's opinion in full while candidly acknowledging that Texas's school finance method produced a system ``that can fairly be described as chaotic and unjust'' (\href{https://supreme.justia.com/cases/federal/us/411/1/}{Justia}). Nonetheless, he emphasized that the Equal Protection Clause ``confers no substantive rights and creates no substantive liberties'' (noting voting equality as a narrow exception) and instead merely polices invidious legislative classifications (\href{https://supreme.justia.com/cases/federal/us/411/1/}{Justia}). Because Texas's system neither created objectively identifiable suspect classes nor impinged on an independent constitutional right, and rested on grounds rationally related to local control, invalidating it would represent ``an extraordinary departure from principled adjudication'' (\href{https://supreme.justia.com/cases/federal/us/411/1/}{Justia}).

\subsection{Dissenting Opinion (Justice Brennan)}

Justice Brennan joined Justice White's dissent finding the Texas system devoid of any rational basis, and wrote separately to reject the majority's ``rather distressing assertion'' that a right is fundamental only if explicitly or implicitly guaranteed by the constitutional text (\href{https://supreme.justia.com/cases/federal/us/411/1/}{Justia}). Endorsing Justice Marshall's nexus principle, Brennan argued that because education is ``inextricably linked'' to voting and First Amendment rights of speech and association, classifications affecting education must receive strict scrutiny---which Texas conceded its system could not survive (\href{https://supreme.justia.com/cases/federal/us/411/1/}{Justia}).

\subsection{Dissenting Opinion (Justice White)}

Justice White, joined by Justices Douglas and Brennan, accepted local control as a legitimate state objective but concluded that Texas's financing scheme failed even \emph{rational basis} review because the means chosen did not rationally advance that goal (\href{https://supreme.justia.com/cases/federal/us/411/1/}{Justia}). While Alamo Heights enjoyed a genuine choice to enrich its schools, Edgewood was practically and legally precluded from doing so: to match the highest Bexar County yield, Edgewood would have needed a prohibitive tax rate of \$5.76 per \$100 valuation, yet Texas law capped maintenance taxes at \$1.50 per \$100 (\href{https://supreme.justia.com/cases/federal/us/411/1/}{Justia}). White argued that extending a meaningful option to property-rich districts while denying it to property-poor districts rendered equal protection ``no more than an empty gesture,'' and that affirming the District Court would still leave Texas free to adopt any rational, fiscally neutral system (\href{https://supreme.justia.com/cases/federal/us/411/1/}{Justia}).

\subsection{Dissenting Opinion (Justice Marshall)}

Justice Marshall, joined by Justice Douglas, delivered a sweeping dissent condemning the majority's decision as ``a retreat from our historic commitment to equality of educational opportunity'' (\href{https://supreme.justia.com/cases/federal/us/411/1/}{Justia}). Marshaling trial evidence of stark resource gaps---100\% of Alamo Heights teachers held college degrees compared to 80.02\% in Edgewood, where 47\% taught on emergency permits---Marshall argued that the Equal Protection Clause forbids unjustifiable inequality in educational opportunity, not merely gross inadequacy (\href{https://supreme.justia.com/cases/federal/us/411/1/}{Justia}).

Doctrinally, Marshall attacked the majority's ``rigidified'' two-tier framework (strict scrutiny versus mere rationality) as inconsistent with precedent, advocating instead a ``spectrum of standards'' (sliding-scale review) calibrated to ``the character of the classification in question, the relative importance to individuals in the class discriminated against of the governmental benefits that they do not receive, and the asserted state interests'' (\href{https://supreme.justia.com/cases/federal/us/411/1/}{Justia}). Pointing to procreation (\textit{Skinner v. Oklahoma}), state voting (\textit{Reynolds v. Sims}), and criminal appeals (\textit{Griffin v. Illinois}), he argued that fundamental interests need not be independently enumerated in the Constitution; rather, as the nexus between a nonconstitutional interest and explicit constitutional guarantees (such as speech and political participation) draws closer, scrutiny must intensify (\href{https://supreme.justia.com/cases/federal/us/411/1/}{Justia}). Finally, Marshall dismissed Texas's local-control rationale as a ``mere sham,'' noting that Texas already regulated minute details of public education and that fiscal-neutrality alternatives (such as district power equalizing) could preserve local governance without tying a child's education to neighborhood property wealth (\href{https://supreme.justia.com/cases/federal/us/411/1/}{Justia}).

\section{Examples: Future Applications}

\begin{enumerate}
\item \textbf{Same-Side (Constitutional under \textit{Rodriguez}):} A state authorizes non-reorganized rural school districts to charge families a modest annual user fee (e.g., \$97 per child) for public school bus transportation, without waiving the fee for indigent families, where the child still attends school by arranging private transport. \emph{Why:} Modeled on \href{https://supreme.justia.com/cases/federal/us/487/450/}{\textit{Kadrmas v. Dickinson Public Schools}}, 487 U.S. 450 (1988), this involves neither a suspect class nor an absolute denial of education, so \textit{Rodriguez} mandates deferential rational basis review, which the state satisfies by conserving funds and encouraging district reorganization (\href{https://supreme.justia.com/cases/federal/us/487/450/}{Justia}).
\item \textbf{Same-Side (Constitutional under \textit{Rodriguez}):} A state funds public libraries and municipal parks exclusively through local municipal property taxes, causing wealthy suburbs to operate state-of-the-art libraries while low-wealth cities operate aging branches with limited hours. \emph{Why:} Like public education in \textit{Rodriguez}, libraries and parks are not explicitly or implicitly guaranteed by the federal Constitution, and inter-municipal property-wealth disparities are rationally related to local fiscal autonomy (\href{https://supreme.justia.com/cases/federal/us/411/1/}{Justia}).
\item \textbf{Opposite-Side (Unconstitutional):} A state enacts a statute withholding all state funding for the education of undocumented school-age children and authorizing local school districts to bar them completely from public schools unless they pay full tuition. \emph{Why:} Modeled on \href{https://www.law.cornell.edu/supremecourt/text/457/202}{\textit{Plyler v. Doe}}, 457 U.S. 202 (1982), this triggers the \emph{absolute deprivation} caveat Powell left open in \textit{Rodriguez}: completely denying basic education to innocent children creates a permanent underclass and cannot survive heightened rationality review absent a substantial state interest (\href{https://www.law.cornell.edu/wex/plyler_v_doe_\%281982\%29}{Cornell LII Wex}).
\item \textbf{Opposite-Side (Unconstitutional):} A state imposes a \$500 mandatory filing fee on all criminal defendants seeking a first appeal of right, with no indigency waiver, completely barring poor defendants from appellate review. \emph{Why:} As Powell expressly distinguished in \textit{Rodriguez} (citing \textit{Griffin v. Illinois} and \textit{Douglas v. California}), this imposes an absolute wealth-based deprivation on a definable class of indigent persons accessing a core criminal-justice safeguard (\href{https://supreme.justia.com/cases/federal/us/411/1/}{Justia}).
\item \textbf{Fence-Sitter (Close Case):} Students in an urban public school district with collapsed roofs, severe teacher vacancies, no textbooks, and a 5\% literacy proficiency rate sue the state under the Fourteenth Amendment, alleging that state underfunding has deprived them of ``access to literacy''---the minimal basic skills necessary to exercise First Amendment and voting rights. \emph{Why:} This mirrors modern litigation such as \textit{Gary B. v. Whitmer} (6th Cir. 2020, vacated upon settlement) (\href{https://en.wikipedia.org/wiki/San_Antonio_Independent_School_District_v._Rodriguez}{Wikipedia}; \href{https://houstonlawreview.org/article/18015-reversing-_rodriguez_-a-siren-call-to-a-dangerous-shoal}{Houston L. Rev.}). Plaintiffs invoke Powell's explicit reservation in \textit{Rodriguez} regarding a ``basic minimal'' quantum of education needed for speech and voting, while the state invokes \textit{Rodriguez}'s broader holding against judicially created positive rights and federal oversight of state educational quality (\href{https://supreme.justia.com/cases/federal/us/411/1/}{Justia}).
\end{enumerate}

\section{Critique}

\textbf{Originalist, Textualist, and Federalist Perspectives:} Conservative, originalist, and federalist scholars have generally defended \textit{Rodriguez}'s core restraint. Judge Jeffrey S. Sutton, writing in the \textit{Virginia Law Review}, argued that Powell's refusal to constitutionalize school finance under the federal Fourteenth Amendment wisely preserved federalism and catalyzed a vibrant wave of state-court litigation and legislative reform under explicit state constitutional education clauses---including the Texas Supreme Court's own \textit{Edgewood} decisions (\href{https://www.virginialawreview.org/wp-content/uploads/2020/12/1963-1.pdf}{Sutton, Va. L. Rev.}). Similarly, commentators warning against overruling \textit{Rodriguez} emphasize that federal courts lack judicially manageable standards to define ``adequate'' funding or resolve empirical debates over school spending (\href{https://houstonlawreview.org/article/18015-reversing-_rodriguez_-a-siren-call-to-a-dangerous-shoal}{Houston L. Rev.}). Interestingly, however, a newer school of \emph{originalist} scholarship---led by Professor Derek W. Black in the \textit{Notre Dame Law Review}---challenges \textit{Rodriguez} on originalist grounds, arguing that the Reconstruction Congress required states to guarantee public education in their state constitutions as a condition of readmission and statehood, linking access to basic education directly to Fourteenth Amendment citizenship (\href{https://ndlawreview.org/wp-content/uploads/2025/08/The-Fundamental-Right-to-Education.pdf}{Black, Notre Dame L. Rev.}).

\textbf{Progressive and Doctrinal Perspectives:} Progressive legal scholars have widely condemned \textit{Rodriguez} as one of the most consequential and damaging decisions of the post-\textit{Brown} era. Critics such as Erwin Chemerinsky and critical race scholars argue that by refusing to recognize poverty as a suspect classification or education as a fundamental right---and by treating race and wealth as analytically separate---the Court entrenched ``separate and unequal'' public schools across metropolitan boundary lines (\href{https://en.wikipedia.org/wiki/San_Antonio_Independent_School_District_v._Rodriguez}{Wikipedia}; \href{https://lawcat.berkeley.edu/record/1124439/files/fulltext.pdf}{Berkeley Law Scholarship}; \href{https://loyola-chicago-law-journal.scholasticahq.com/article/92496-_san-antonio-independent-school-district-v-rodriguez_-at-fifty-contingencies-consequences-and-calls-to-action/attachment/193053.pdf}{Loyola Chi. L.J.}). Doctrinally, scholars note that Powell's rigid two-tier Equal Protection model proved unstable even on its own terms, as the Court soon adopted intermediate scrutiny for sex and nonmarital children and applied heightened rationality in \href{https://www.law.cornell.edu/supremecourt/text/457/202}{\textit{Plyler v. Doe}} (\href{https://www.law.cornell.edu/constitution-conan/amendment-14/section-1/educational-opportunity}{Constitution Annotated}).

\textbf{Underlying Competing Values:} \textit{Rodriguez} sits at the crossroads of \emph{substantive equality of opportunity}---ensuring that a child's life prospects do not depend on the accident of neighborhood property wealth---and \emph{democratic federalism and judicial restraint}, which reserve complex decisions over taxation, resource allocation, and local school governance to state legislatures and state constitutions rather than life-tenured federal judges (\href{https://supreme.justia.com/cases/federal/us/411/1/}{Justia}; \href{https://www.virginialawreview.org/wp-content/uploads/2020/12/1963-1.pdf}{Sutton, Va. L. Rev.}).

\section{Key Quotations}

\begin{quote}
``The system of alleged discrimination and the class it defines have none of the traditional indicia of suspectness: the class is not saddled with such disabilities, or subjected to such a history of purposeful unequal treatment, or relegated to such a position of political powerlessness as to command extraordinary protection from the majoritarian political process.'' \\
--- Justice Powell, majority opinion (\href{https://supreme.justia.com/cases/federal/us/411/1/}{411 U.S. at 28, Justia})
\end{quote}

\begin{quote}
``It is not the province of this Court to create substantive constitutional rights in the name of guaranteeing equal protection of the laws. Thus, the key to discovering whether education is `fundamental' is not to be found in comparisons of the relative societal significance of education, as opposed to subsistence or housing. Nor is it to be found by weighing whether education is as important as the right to travel. Rather, the answer lies in assessing whether there is a right to education explicitly or implicitly guaranteed by the Constitution.'' \\
--- Justice Powell, majority opinion (\href{https://supreme.justia.com/cases/federal/us/411/1/}{411 U.S. at 33--34, Justia})
\end{quote}

\begin{quote}
``Education, of course, is not among the rights afforded explicit protection under our Federal Constitution. Nor do we find any basis for saying it is implicitly so protected.'' \\
--- Justice Powell, majority opinion (\href{https://supreme.justia.com/cases/federal/us/411/1/}{411 U.S. at 35, Justia})
\end{quote}

\begin{quote}
``The method of financing public schools in Texas, as in almost every other State, has resulted in a system of public education that can fairly be described as chaotic and unjust. It does not follow, however, and I cannot find, that this system violates the Constitution of the United States.'' \\
--- Justice Stewart, concurring (\href{https://supreme.justia.com/cases/federal/us/411/1/}{411 U.S. at 59, Justia})
\end{quote}

\begin{quote}
``The Court apparently seeks to establish today that equal protection cases fall into one of two neat categories which dictate the appropriate standard of review---strict scrutiny or mere rationality. But this Court's decisions in the field of equal protection defy such easy categorization. A principled reading of what this Court has done reveals that it has applied a spectrum of standards in reviewing discrimination allegedly violative of the Equal Protection Clause.'' \\
--- Justice Marshall, dissenting (\href{https://supreme.justia.com/cases/federal/us/411/1/}{411 U.S. at 98--99, Justia})
\end{quote}

\end{document}
